% Physical pages 3 and 4 of the concise study: the texts of the Mass in their
% order, what each says, keeps and omits, with the few points of the Fathers a
% reader needs to follow them. The developed comparison of the three
% interpretations begins on the fifth page.

\section{The Propers: Themes and Movement}
\zlabel{triptych:concise:themes:start}

On the First Sunday of Advent the prophet asks God to tear the sky open and
come down; the Gospel commands servants to stay awake for a master who has
gone away and will come back; the Preface gives thanks for two comings of Christ, the first in lowly
flesh and the second in glory. In Advent the readings are not taken in
course, as on the Sundays of the year, but chosen for the season: on this
Sunday a Gospel of the Lord's coming at the end of time, a prophecy of the
Messiah from Isaiah, and an apostle's exhortation (\work{Praenotanda} 67,
93). The readings are Year~B's; the antiphons, prayers and Preface are the
Missal's for this Sunday in all three years.

The Mass opens with Psalm 25 (24): \latin{Ad te levavi animam meam}, ``To thee
have I lifted up my soul.'' The speaker trusts in God and asks not to be
shamed or mocked by his enemies, ``for none of them that wait on thee shall be
confounded.'' At that last phrase the antiphon does not follow the Clementine
Vulgate or the \work{Nova Vulgata}, which read \latin{qui sustinent te}; it
reads \latin{qui te exspectant}, ``who wait for thee''. The verb
\latin{exspectare} comes again in the Latin of the Mass: Paul's Corinthians are ``waiting for the manifestation
of our Lord'', and the Preface ends with those who dare to wait. The Collect
then asks God for a will: that he give his faithful the will
(\latin{voluntatem}) to go out with just works to meet his Christ as he comes
(\latin{Christo tuo venienti \ldots\ occurrentes}), so that, set at his right
hand, they may possess the kingdom of heaven. Christ is spoken of in the
present participle, as one already on the road, and the right hand is where
the Son of man sets the sheep in Matthew's judgment scene (Mt 25:33).

The first reading takes three pieces of a long prayer in Isaiah 63--64 that
is spoken among ruins. The prayer begins by recalling the Lord's mercies at
the Exodus, when ``the angel of his presence saved them'' and Moses led them
through the deep, and asks where that Lord is now (63:9--15). The people say
that their enemies have trodden down the sanctuary, and the prayer ends before
the house ``where our fathers praised thee'', burnt with fire (63:18; 64:10,
Douay 64:11). In the reading they call the Lord ``our father, our
redeemer'' and then ask, ``Why hast thou made us to err, O Lord, from thy
ways: why hast thou hardened our heart, that we should not fear thee?''
(63:17). They ask him to return for the sake of his servants, and cry, ``O
that thou wouldst rend the heavens, and wouldst come down: the mountains would
melt away at thy presence'' (63:19b). The Latin \work{Ordo lectionum Missae}
sets this cry over the reading as its title. Two verses on, the Latin of the reading has the people say \latin{Descendisti}, ``thou didst come down'', and
the mountains melt a second time, now as a thing done. The English proclaimed
in the United States has no words for \latin{Descendisti}. Jerome, who
translated the prophet from the Hebrew, took that past tense for the sign of a
prayer heard. The next verse says, ``the eye hath not seen, O God, besides
thee, what things thou hast prepared for them that wait for thee'' (64:3;
Douay 64:4). Paul writes ``That eye hath not seen, nor ear heard'' in this
same letter to Corinth (1~Cor 2:9), and Jerome insists that Paul is giving
Isaiah's sense there, not quoting an apocryphal book. The confession
follows. ``Thou hast met him that rejoiceth, and
doth justice'': the Latin verb, \latin{occurristi}, is the Collect's, so that
in Isaiah God goes out to meet the just and in the Collect the faithful ask to
go out to meet Christ. Then, ``behold thou art angry, and we have sinned'';
``all our justices as the rag of a menstruous woman: and we have all fallen
as a leaf, and our iniquities, like the wind, have taken us away''; ``there is
none that
calleth upon thy name''; ``thou hast hid thy face from us''. The reading ends,
``And now, O Lord, thou art our father, and we are clay: and thou art our
maker, and we all are the works of thy hands'' (64:4--7; Douay 64:5--8).

The psalm, Psalm 80 (79), is a lament for a defeated people, addressed to the
Shepherd of Israel enthroned on the cherubim, the winged figures over the ark.
Its middle, which is not sung, tells Israel's history as a vine brought out of
Egypt and planted, its walls now broken and the boar out of the wood laying it
waste. The strophes that are sung ask God to stir up his might and come, to
look down from heaven and visit the vine, and to lay his hand ``upon the man
of thy right hand: and upon the son of man whom thou hast confirmed for
thyself''; the son of man closes v.~16 as well as v.~18. Then
the people promise, ``we depart not from thee''. The response
is the psalm's own refrain: ``Convert us, O God: and shew us thy face, and we
shall be saved.'' It comes three times in the whole psalm (vv.~4, 8, 20), each
time with a fuller name for God: ``O God'', ``O God of hosts'', ``O Lord God
of hosts''. Augustine began his sermon on the psalm by saying what it
is about: ``The song here is of the Advent of the Lord and of our Saviour
Jesus Christ, and of His vineyard'', Christ and his people, ``Head and Body''
(\work{Enarrationes in Psalmos} 79.1).

Paul's first letter to Corinth opens with thanks. The Corinthians are rich ``in
all utterance and in all knowledge'', and nothing is wanting to them in any
grace as they wait ``for the manifestation of our Lord Jesus Christ''; Christ
``will confirm you unto the end without crime, in the day of the coming of our
Lord''; and ``God is faithful: by whom you are called unto the fellowship of
his Son'' (1~Cor 1:5--9). ``Without crime'' renders the Latin \latin{sine
crimine}, which stands for a Greek word for one against whom no charge can be
brought. The rest of the letter rebukes this church for
factions, for incest, for lawsuits and for disorder at the Lord's Supper, and
calls its members carnal (3:1). Over the reading the \work{Ordo} puts Paul's
participle into the mouth of the assembly: \latin{Revelationem exspectamus},
``We await the revelation.''

The Alleluia verse comes from Psalm 85 (84), which begins by thanking God for
turning away the captivity of Jacob and then asks for mercy again: ``Shew us,
O Lord, thy mercy; and grant us thy salvation'' (v.~8). Augustine, reaching
the verse in a sermon, began with \latin{Hoc est quod cantavimus}, ``This is
what we have sung'': his congregation had sung it before he explained it
(84.9). In the Mass it greets the Gospel.

The Gospel is the close of the last speech Jesus makes in Mark before the
Passion; the next verse is ``Now the feast of the pasch and of the Azymes was
after two days'' (14:1). On the Mount of Olives, opposite the Temple, Peter, James, John and
Andrew have asked him when its stones will fall. He has spoken of deceivers
and wars, of the darkened sun and the Son of man coming in the clouds, and has
then refused them the date: ``of that day or hour no man knoweth, neither the
angels in heaven, nor the Son, but the Father'' (13:32). That verse is not
read. The reading opens with two commands, \latin{Videte, vigilate}, ``look,
keep awake''; \latin{vigilare} is to stay awake through the night. A man going abroad leaves his house, gives his servants
authority over every work, and commands the doorkeeper to watch. He may come
at evening, at midnight, at cockcrow or in the morning, the four watches of
the Roman night, and he does not come back within the story. The last verse
reaches past the four men: ``And what I say to you, I say to all: Watch.''
The \work{Ordo}'s title over the Gospel is \latin{Vigilate: nescitis quando
dominus domus veniat}, ``Watch: you know not when the lord of the house
comes.''

The Prayer over the Offerings asks God to accept gifts offered from what he
has given (\latin{de tuis \ldots\ collata beneficiis}), and asks that the
devotion he grants in this passing time become the reward of eternal
redemption. The Preface gives thanks for Christ's two comings: the first in
the lowliness of the flesh he took (\latin{in humilitate carnis assumptae}),
which opened the way of salvation; the second in the glory of his majesty,
when we who now keep watch and dare to wait (\latin{vigilantes nunc audemus
exspectare}) will receive openly what was promised. The Communion Antiphon is
the verse of Psalm 85 (84) five verses after the Alleluia's: ``the Lord will
give goodness: and our earth shall yield her fruit''. In the psalm it follows
the meeting of mercy and truth, justice and peace, when ``Truth is sprung out
of the earth'' (vv.~11--12). The Prayer after
Communion asks that the mysteries profit those who receive them, and gives
the reason: through them God already teaches his people, while they live among
things that pass away, to love what is heavenly and to cling to what lasts.
The solemn blessing, where it is used, asks eternal life, when the Son comes
again in majesty, for those who rejoice that he came in the flesh.

Three questions rise from these texts. Has the prophet's cry been answered,
and if it has, why does the Church still pray it? Why does Christ hide the
hour, and to whom is ``Watch'' spoken, when most who hear it will die before
the end? And on what ground can a people that confesses it has strayed hope to
be found blameless on the day?
\zlabel{triptych:concise:themes:end}
